\chapter{シグナル}
\label{ch:signal}

\section{シグナルとは}

クオンツ運用で\textbf{シグナル}（アルファ・シグナル、ファクター、特徴量）とは、各時点 $t$・各資産 $i$ について計算される数 $s_{i,t}$ で、将来の（相対的な）リターンを予測するものをいう。

\begin{itemize}
\item 時点 $t$ で\textbf{実際に知られていた}情報だけから計算する。
\item 単独の銘柄の上げ下げより、\textbf{銘柄の間の比較}（どれが相対的に上がるか）を予測するものが多い（横断面のシグナル）。
\item 一つ一つは弱い（IC が 0.01〜0.05 程度）。多くを組み合わせて使う。
\end{itemize}

先物やマクロの戦略では、一つの資産の時系列の上げ下げを予測するシグナル（時系列のシグナル）も多い。

\section{シグナルの種類}

\begin{table}[H]
\centering\small
\caption{代表的なシグナル（株式の場合）}
\label{tab:signals}
\begin{tabular}{@{}P{0.13\linewidth}P{0.3\linewidth}P{0.3\linewidth}P{0.18\linewidth}@{}}
\toprule
分類 & 例 & 計算の例 & 効く期間 \\
\midrule
価格 & モメンタム & 過去 12 か月から直近 1 か月を除いたリターン（12-1） & 数か月〜1 年 \\
 & 短期リバーサル & 直近 1 か月（または 1 週）のリターンの逆 & 数日〜1 か月 \\
 & 低ボラティリティ & 過去 1 年の日次リターンの標準偏差の逆、ベータの逆 & 数か月〜年 \\
\midrule
財務 & 割安 & 簿価時価比率、利益÷株価、キャッシュフロー÷企業価値 & 年 \\
 & 収益性・クオリティ & 売上総利益÷総資産、ROE、利益の安定性、負債の少なさ & 年 \\
 & 投資・発行 & 総資産の伸び、株式の発行（伸びが大きいほど低いリターン） & 年 \\
 & 会計の質 & 発生主義の利益の割合（アクルーアル）の大きさの逆 & 年 \\
\midrule
予想・発表 & 決算の驚き & SUE＝（実際の利益 − 予想）÷ 標準偏差 & 数日〜数か月 \\
 & 予想の修正 & アナリストの利益予想の上方修正の数 − 下方修正の数 & 数週〜数か月 \\
\midrule
需給 & 空売りの残高 & 空売りの株数 ÷ 浮動株（多いほど低いリターン） & 数週〜数か月 \\
 & 資金の流れ & 投資信託の資金の流出入による売買の圧力 & 数週〜数か月 \\
\midrule
代替データ & テキスト & 決算発表文・ニュース・SNS の語調、LLM の読み取り & 数日〜数週 \\
 & その他 & カード決済、衛星画像、ウェブの閲覧、求人、アプリの利用 & 数週〜四半期 \\
\bottomrule
\end{tabular}
\end{table}

先物・マクロでは、次のものが代表的である。

\begin{itemize}
\item \textbf{時系列モメンタム（トレンド）}：その資産自身の過去 1〜12 か月のリターンの符号、移動平均の交差。
\item \textbf{キャリー}：価格が変わらなかったときに得られる収益。通貨なら金利の差、債券なら利回りと年限の効果、商品なら先物の期近と期先の価格差。
\item \textbf{バリュー}：購買力平価からの乖離（通貨）、実質の利回り（債券）、長期の平均からの乖離。
\end{itemize}

\section{シグナルの作り方}

生のデータから使えるシグナルを作るまでに、次の処理をする。

\subsection{時点の正しいデータ}

その時点で知られていた値だけを使う。

\begin{itemize}
\item \textbf{発表の遅れ}：12 月期の決算の数値が公表されるのは、翌年の 2〜3 月である。「12 月末の簿価」を 12 月末の売買に使うと、未来の情報を使ったことになる（先読み、look-ahead）。Fama・French は、前年度の財務データを 6 月末から使う。
\item \textbf{訂正}：財務データは後から訂正される。データベースの「最新の値」は、当時の値とは違う。当時の値を残したデータ（ポイント・イン・タイムのデータ）が要る。
\item \textbf{銘柄の範囲}：今日存在する銘柄だけで過去を検証すると、倒産・上場廃止した銘柄が抜け、成績が良く見える（生存者バイアス）。
\end{itemize}

\subsection{横断面の標準化}

シグナルの水準は、時期によって変わる（例えば、全体の利益の水準）。そこで、毎時点、銘柄の間で標準化する。
\begin{equation}
z_{i,t} = \frac{s_{i,t} - \bar{s}_t}{\mathrm{sd}_t(s)} \qquad \text{または} \qquad
z_{i,t} = \Phi^{-1}\!\left(\frac{\mathrm{rank}_t(s_{i,t}) - 0.5}{N_t}\right)
\end{equation}
後者は順位を正規分布の分位点に直したもの（正規スコア）で、外れ値の影響を受けにくい。LLM-BLM の BLM の見通しでも、点数を正規スコアにした。

\subsection{外れ値の処理}

上下 1\% や平均から 3 標準偏差の値で丸める（ウィンソライズ）。財務の比率は分母が小さいと極端な値になるので、特に必要である。

\subsection{中立化}

シグナルが、意図しないファクターを含んでいることが多い。例えば「割安」は特定の業種（金融、エネルギー）に偏り、「低ボラティリティ」は大型株と公益事業に偏る。毎時点、シグナルを業種のダミーや規模・ベータに回帰し、その残差をシグナルとする（中立化）。
\begin{equation}
z_{i,t} = \sum_{g} c_{g,t}\,\mathbb{1}[i \in g] + d_t\,\ln(\mathrm{size}_{i,t}) + \tilde{z}_{i,t}
\end{equation}
残差 $\tilde{z}_{i,t}$ は、業種と規模の中で比べたシグナルになる。

\section{シグナルの評価}

\subsection{IC}

時点 $t$ のシグナル $z_{i,t}$ と、その後のリターン $r_{i,t+h}$ の銘柄の間の相関を\textbf{情報係数}（IC）と呼ぶ。外れ値に強い順位相関（スピアマン）がよく使われる（rank IC）。
\begin{equation}
\IC_t = \mathrm{corr}_i\big(\mathrm{rank}(z_{i,t}),\ \mathrm{rank}(r_{i,t+h})\big)
\end{equation}
時点ごとの $\IC_t$ の平均と、その安定性（平均÷標準偏差、IC の IR）を見る。

\subsection{分位のポートフォリオ}

シグナルで銘柄を 5 つまたは 10 の組に分け、組ごとの平均リターンを見る。平均リターンが組の順に単調に並ぶか、上位と下位の差（スプレッド）が大きいか、上位と下位のどちらに収益が偏っているかを確かめる。空売りが難しい銘柄（小型・流動性が低い）で下位の収益が大きいなら、実際にはとれない。

\subsection{減衰}

シグナルを計算した後、$h = 1, 2, \dots$ 日のリターンに対する IC を描くと、シグナルの情報がどれだけ速く価格に入るかが分かる（減衰の曲線）。

\begin{figure}[H]
\centering
\begin{tikzpicture}[x=9mm,y=28mm]
\draw[arr] (0,0)--(11,0) node[below left,font=\footnotesize]{シグナルを計算してからの日数};
\draw[arr] (0,0)--(0,1.15) node[above,font=\footnotesize]{累積の IC（相対値）};
\draw[c2,line width=1pt,domain=0:10.5,samples=60] plot (\x,{1-exp(-3*\x)});
\draw[c1,line width=1pt,domain=0:10.5,samples=60] plot (\x,{1-exp(-0.25*\x)});
\node[lab,c2,anchor=west] at (1.2,1.05) {速い減衰（発表の当日に入りきる）};
\node[lab,c1,anchor=west] at (5.5,0.62) {遅い減衰（数週間かけて入る）};
\foreach \x in {0,2,4,6,8,10} \node[lab,below] at (\x,0) {\x};
\end{tikzpicture}
\caption{シグナルの減衰の概念図。速く減衰するシグナルは、すぐに売買しなければ価値がなく、売買の回数も多くなる。}
\label{fig:decay}
\end{figure}

\subsection{回転率と容量}

シグナルの順位が毎期どれだけ入れ替わるか（自己相関）で、売買の量が決まる。シグナルの自己相関が低いほど、回転率が高く、取引コストがかかる。同じ IC でも、ゆっくり変わるシグナルのほうが、実際の運用では価値が高いことが多い。

\subsection{ほかのシグナルとの相関}

新しいシグナルが、既知のシグナル（モメンタム、割安など）と強く相関しているなら、新しい情報は少ない。既知のシグナルに回帰した残差の IC（増分の IC）を見る。

\section{シグナルの組み合わせ}

多くのシグナル $z^{(1)}, \dots, z^{(K)}$ を一つの予測にまとめる方法は、単純なものから順に次のとおり。

\begin{enumerate}
\item \textbf{等しい重み}：標準化したシグナルの平均。重みを推定しないので、過学習しない。意外に強い。
\item \textbf{IC で重み}：過去の IC（あるいは IC の IR）に比例した重み。
\item \textbf{回帰}：過去の将来リターンを目的変数に、シグナルを説明変数にした回帰の係数を重みにする（リッジなどで縮める）。シグナルの間の相関を考慮できる。
\item \textbf{機械学習}：決定木やニューラルネットで、非線形な組み合わせや相互作用を学ぶ（\chapref{ch:model}）。
\end{enumerate}

重みの推定には誤差があり、推定の期間で良かったシグナルが、次の期間も良いとは限らない。重みを等しい重みに向けて縮めること、重みの変化をゆっくりにすることが、実務の定石である。

\section{シグナルの寿命}

シグナルは、時間とともに弱くなることが多い。

\begin{itemize}
\item \textbf{公表}：論文で公表されたシグナルは、公表の後に収益が小さくなる（第 II 部）。
\item \textbf{混雑}：同じシグナルで売買する資金が増えると、収益が小さくなり、一斉に売るときの損が大きくなる。
\item \textbf{構造の変化}：データの入手が容易になる（代替データの普及）、取引の仕組みが変わる（手数料の自由化、ティックの縮小）。
\end{itemize}

運用会社は、シグナルの成績を常に監視し、新しいシグナルを探し続ける。研究の部門の大きな仕事は、「まだ広く使われていないデータ」と「既知のデータの新しい使い方」を探すことである。

\begin{projbox}
LLM の決算発表文の読み取りは、表~\ref{tab:signals} の「テキスト」の一つである。LLM-BLM の実験では、2 日の反応に対する IC は高かったが、減衰の曲線（\figref{fig:decay}）でいえば「速い減衰」で、発表の翌日以降のリターンへの予測力はほとんど残らなかった。この種のシグナルを運用で使うには、発表の直後（数分〜当日）に売買するか、反応の遅い銘柄（注目の低い小型株：実験 S1 の仮説）を探す必要がある。
\end{projbox}
